\documentclass[10pt,a4paper]{article}
\usepackage{amsmath,amssymb,amsthm,graphicx}
\usepackage{amsmath,amsthm,amssymb,amsfonts}
\usepackage[shortlabels]{enumitem}
\newcommand{\T}{\mathcal{T}}
\newcommand{\D}{\mathcal{D}}
\newcommand{\CC}{\mathcal{C}}
\newcommand{\oo}{\mathcal{O}}
\newcommand{\0}{\varnothing}
\newcommand{\Z}{\mathbb{Z}}
\newcommand{\N}{\mathbb{N}} 
\newcommand{\Q}{\mathbb{Q}} 
\newcommand{\R}{\mathbb{R}}
\newcommand{\F}{\mathcal{F}}
\newcommand{\p}{\mathfrak{p}}
\newcommand{\q}{\mathfrak{q}}
\newcommand{\m}{\mathfrak{m}}
\newcommand{\ann}{\textnormal{ann}}
\newcommand{\hm}{\textnormal{Hom}}
\newcommand{\id}{\textnormal{id}}
\newcommand{\ind}{\textnormal{Ind}}
\newcommand{\im}{\textnormal{im }}
\newcommand{\asdim}{\textnormal{asdim}}
\newcommand{\asi}{\textnormal{asi}}
\newcommand{\inv}{^{-1}}
\newcommand{\st}{\ | \ }
\newcommand{\tor}{\textnormal{Tor}}
\newcommand{\len}{\mathcal{l}}
\newcommand{\gal}{\textnormal{Gal}}
\newcommand{\leg}[2]{\left( \frac{#1}{#2} \right)_n}
\newcommand{\sleg}[2]{\left( #1/#2 \right)_n}
\newcommand{\pp}{(\oo/\p)^\times}
\newcommand{\qq}{(\oo/\q)^\times}
\newcommand{\tr}{\textnormal{Tr}}
\newcommand{\cay}{\textnormal{Cay}}
\newcommand{\PMEU}{PME\textsubscript{U}}
\newcommand{\res}{\upharpoonright}
\newtheorem{theorem}{Theorem}
\newtheorem{corollary}{Corollary}
\newtheorem{lemma}{Lemma} 
\newtheorem{prop}{Proposition}  
\newtheorem{prob}{Problem} 
\title{Borel degree-plus-one edge coloring from asymptotic separation index one}
\author{Felix Weilacher}
\date{arxiv:2104.14646v2; CC BY 4.0}
\begin{document}
\maketitle
\noindent\textit{Editorial note.} Only the ASI1 degree-plus-one theorem is counted. The two ray/third-color lemmas are proved in full as solution obligations. No theorem establishing ASI1 for particular groups is assumed or counted.\par\medskip
\noindent\textit{Editorial note.} For any ray, the argument applies to every tail, so a point in U0 may be followed by a later point outside its2d-neighborhood when checking the annulus crossings184.\par\medskip
\noindent\textit{Editorial note.} The third color used for different generator pairs occurs only on edges whose endpoints lie in their pairwise disjoint V\_i, which is the reason it creates no cross-generator color conflict.\par\medskip
\noindent\textit{Editorial scope.} One Borel edge-coloring theorem with its complete local geometric/Borel construction. Borel maximal independent sets and finite-component choice are the explicit standard inputs. The finite-extension consequences and rank-specific abelian classification are outside this counted task.\par\medskip
\noindent\textit{Allowed standard prerequisites (editorial).}\begin{itemize}
\item Borel maximal independent sets in locally finite Borel graphs. Every locally finite Borel graph admits a Borel maximal independent set.. Kechris-Solecki-Todorcevic1999, Proposition4.2, explicitly invoked at193.. Applied to a locally finite line graph to obtain a maximal matching inside each selected ray-hitting region.
\item Borel finite choices. Finite connected components and finitely many colorings per finite component permit a Borel choice, using a fixed Borel order.. Choose alternating two-edge-colorings of finite paths or even cycles; the global disjoint regions and edge cuts are not given.
\end{itemize}\medskip
Let $X$ be a standard Borel space and $G \subset X \times X$ a Borel graph on $X$. An \textit{edge coloring} of $G$ is a function $c:G \rightarrow Y$ which sends adjacent edges in $G$, that is, edges sharing a vertex, to distinct elements of $Y$, or ``colors''. Such a $c$ is called a \textit{$k$-coloring} if $|Y| = k$. The \textit{edge chromatic number} of $G$, denoted $\chi'(G)$, is the smallest cardinal $k$ such that $G$ admits a $k$-coloring. The \textit{Borel edge-chromatic number} of $G$, denoted $\chi_B'(G)$, is the smallest cardinal $k$ such that $G$ admits a Borel $k$-coloring. See \cite{KM} for a survey of these and related notions.

The Borel graphs of interest to us in this paper will be those generated by Borel actions of finitely generated groups. A \textit{marked group} is a pair $(\Gamma,S)$ where $\Gamma$ is a group and $S$ is a finite symmetric generating set for $\Gamma$ not containing 1. We will sometimes omit the $S$ if it will not cause confusion. Let $X$ be a standard Borel space and $a:\Gamma \curvearrowright X$ a free Borel action. Let $G(a,S)$ denote the \textit{Schreier graph} generated by this action. This is the graph defined by setting $x,y \in X$ adjacent if and only if $\gamma \cdot_a x = y$ for some $\gamma \in S$.

Note that all the notions in the above paragraph make sense even if $S$ does not generate the entire group $\Gamma$. We shall sometimes use them in this case.

We start with the discrete versions of these notions: Let $G$ be a locally finite graph on a set $X$, and $N$ a positive integer. Let $G^{\leq N}$ denote the \textit{distance $N$-graph} of $G$. This is the graph which sets two vertices adjacent if and only if there is a path from one to the other of length at most $N$.

The \textit{asymptotic separation index} of $G$, denoted $\asi(G)$ is the infimum over $s \in \omega$ such that the following holds: For every positive integer $N$, $X$ can be partitioned into $s+1$ sets, say $U_0,\ldots,U_s$, such that for each $0 \leq i \leq s$, the induced graph $G^{\leq N} \res U_i$ has only finite connected components. The \textit{asymptotic dimension} of $G$, denoted $\asdim(G)$, is defined similarly, except now the connected components of each $G^{\leq N} \res U_i$ are required to have a uniform (finite) bound on their diameter. 

 Let $G$ now be a Borel graph on a standard Borel space $X$. The \textit{Borel asymptotic separation index} and \textit{Borel asymptotic dimension} of $G$, denoted $\asi_B(G)$ and $\asdim_B(G)$ respectively, are defined exactly as their discrete counterparts were, except that now all the $U_i$'s are required to be Borel sets.
 
 Now let $\Gamma$ be a finitely generated group, and $a:\Gamma \curvearrowright X$ a free Borel action. It is easy to verify that none of the above numbers for the graph $G(a,S)$ depend on the choice of generating set $S$ for $\Gamma$. We can thus talk about the (Borel) asymptotic separation index and (Borel) asymptotic dimension of the action $a$ itself.

Let $G$ be a graph on a set $X$. An \textit{injective $G$-ray} is an injective infinite sequence $x_0,x_1,\ldots \in X$ such that $(x_i,x_{i+1}) \in G$ for all $i \in \omega$. Given a set $U \subset X$ and an $n \in \omega$, let us write $B(U,n)$ for the ball of radius $n$ around $U$. That is, the set of all points $x \in X$ whose path distance to $U$ is less than or equal to $n$. For $m \in \omega$, let us write $A(U,n,m) = B(U,m)-B(U,n)$. One can check that if $G$ and $U$ are Borel, these sets are always Borel.

\begin{theorem}\label{th:d+1}
Let $a:\Gamma \curvearrowright X$ be a free Borel action of a marked group $(\Gamma,S)$ on a standard Borel space $X$ with Borel asymptotic separation index 1, and such that none of the elements of $S$ have odd order. Then $\chi_{B}'(G(a,S)) \leq |S|+1$.
\end{theorem}

\begin{lemma}\label{lem:rainbow}
Suppose $G$ is a locally finite Borel graph on a space $X$ with $\asi_B(G) = 1$, and $d \in \omega$. Then we can find pairwise disjoint Borel sets $V_1,\ldots,V_d \subset X$ such that any injective $G$-ray contains infinitely many edges within each $V_i$.
\end{lemma}

\begin{proof}
Since the tail of an injective $G$-ray is still an injective $G$-ray, it suffices to prove this with ``infinitely many edges'' replaced by ``an edge''. 

Let $X = U_0 \sqcup U_1$ be a Borel partition of $X$ witnessing $\asi_B(G) = 1$ for $N = 4d+1$. That is, such that the graphs $G^{\leq N} \res U_0$ and $G^{\leq N} \res U_1$ both have only finite connected components. Set $V_i = A(U_0,2(i-1),2i)$ for each $1 \leq i \leq d$. Also set $W = X \setminus B(U_0,2d) \subset U_1$. As was mentioned above, these are all still Borel. Let us now show these $V_i$'s work.

Let $x = (x_n)$ be an injective $G$-ray. We first show that $x$ contains some point from $W$. Suppose not. Then for each $n$, there is a $y_n \in U_0$ whose path distance to $x_n$ is less than or equal to $2d$. Now for each $n$, going from $y_n$ to $x_n$, then to $x_{n+1}$, then to $y_{n+1}$, produces a path from $y_n$ to $y_{n+1}$ of length at most $2d+1+2d=N$, so all the $y_n$'s are in the same $G^{\leq N} \res U_0$-connected component. By definition of $U_0$, then, the set $\{y_n \mid n \in \omega\}$ is finite. But now $x$ gives infinitely many points within distance $2d$ of this set, contradicting local finiteness.

Since $W \subset U_1$, the above shows in particular that $x$ contains some point from $U_1$. Similarly, $x$ must contain some point from $U_0$. Let $x_n \in W$ and $x_m \in U_0$. We can assume $m < n$. Now the path distance from $x_i$ to $U_0$ changes by at most 1 each time $i$ is incremented by 1. Thus, since $x_m$ has distance 0 from $U_0$ and $x_n$ has distance greater than $2d$ from $U_0$, for each $1 \leq i \leq d$, there must be some $m < j < n$ such that $x_{j}$ and $x_{j+1}$ have distance $2i-1$ and $2i$ respectively from $U_0$. Then $(x_j,x_{j+1})$ is an edge of $x$ contained in $V_i$, as desired.

\end{proof}

\begin{lemma}\label{lem:Z}
Let $a:\Z \curvearrowright X$ be a free Borel action of $\Z$, and $S$ the usual generating set for $\Z$. Suppose $V \subset X$ is a Borel set with the property that every injective $G(a,S)$-ray contains infinitely many edges within $V$. Then there is a Borel 3-edge coloring of $G(a,S)$, say with the colors $1,2$, and $3$, such that the color 3 only occurs on edges contained within $V$.
\end{lemma}

\begin{proof}
A well known theorem from \cite{KST} (Proposition 4.2) states that for locally finite Borel graphs, it is always possible to find Borel maximal independent sets. Apply this to the edge graph of $G(a,S) \res V$ to get a maximal set of edges $A$ contained within $V$ such that no two edges in $A$ share a vertex. Give all the edges in $A$ the color 3.

By maximality and the defining condition for the set $V$, these edges occur infinitely often in both directions along each $\Z$-orbit. Thus, the connected components of $G(a,S) \setminus A$ are all simply finite paths. These can of course be 2-colored with the colors 1 and 2, and since they are all finite, this can be done in a Borel fashion. (For example, by fixing a Borel linear order on the space of 2-edge colored finite subgraphs of $G(a,S)$, and picking the least 2-coloring in this order for each path above.)
\end{proof}

\begin{proof}
By hypothesis, we may write $S = S_0 \sqcup S_1$, where $S_0$ contains all of our generators of infinite order, and $S_1$ contains all of our elements of even (finite) order. For each $\gamma \in S_1$, the connected components of the graph $G(a,\{\gamma^{\pm 1} \})$ are all even (finite) cycles, or just single edges if $\gamma$ has order 2. Thus they may be 2-colored in a Borel fashion as in the previous proof (or just 1-colored if $\gamma$ has order 2). Thus we may Borel $|S_1|$-color the edges in $G(a,S_1)$.

It now suffices to Borel $|S_0|+1$-color the edges in $G(a,S_0)$ using a disjoint set of colors. Let $S_0 = \{\gamma_1^{\pm 1},\ldots,\gamma_d^{\pm 1}\}$, so that $|S_0| = 2d$. We will use the color set $\{1,\ldots,2d+1\}$. Let $V_1,\ldots,V_d \subset X$ be the sets from Lemma \ref{lem:rainbow}. Since $\langle \gamma_i \rangle \cong \Z$ for each $i$, by Lemma \ref{lem:Z} we may for each $i$ color the edges in $G(a,\{ \gamma_i^{\pm 1} \})$ using the colors $2i-1$, $2i$, and $2d+1$ such that the color $2d+1$ only appears on edges contained within $V_i$. Since the $V_i$'s are pairwise disjoint, this does not cause any color conflicts, so we are done.
\end{proof}

\begin{thebibliography}{9}
\bibitem{KM} A.S Kechris and A.S. Marks, Descriptive graph combinatorics, preprint, 2019; posted in \textnormal{http://www.math.caltech.edu/$\sim$kechris/}

\bibitem{KST} A.S. Kechris, S. Solecki, and S. Todorcevic, Borel chromatic numbers, \textit{Adv. Math.}, 141 (1999), 1-44.
\end{thebibliography}
\end{document}
